\documentclass[11pt,a4paper]{article}
\usepackage[margin=2cm]{geometry}
\usepackage{lmodern}
\usepackage[T1]{fontenc}
\usepackage{amsmath}
\usepackage{graphicx}
\usepackage{float}
\usepackage{microtype}
\usepackage{xcolor}
\usepackage[colorlinks=true,linkcolor=blue!55!black]{hyperref}
\graphicspath{{./}}
\setlength{\parskip}{4pt}
\setlength{\parindent}{0pt}

\newcommand{\fig}[1]{\begin{center}\includegraphics[width=0.92\linewidth]{#1}\end{center}}
\newcommand{\entry}[4]{%
  \subsection*{\texttt{\detokenize{#1}}}
  \fig{#1}
  \textbf{What it shows.} #2\par
  \textbf{How it was made.} #3\par
  \textbf{What it supports.} #4\par
  \vspace{4pt}\hrule\vspace{6pt}}

\title{Figure catalogue: docking estimator and controller}
\author{}
\date{}

\begin{document}
\maketitle
\vspace{-2em}
Every figure exists for both instances, \texttt{demo} (tabletop VEX robot) and \texttt{fullscale} (600~kg barrier robot, simulation only), generated from the implemented MATLAB (\texttt{setup\_paths}, then \texttt{make\_ekf\_figures}, \texttt{make\_mpc\_figures}, \texttt{make\_report\_figures}). Conventions: target frame with the origin at the docked chassis position and $x$ into the target; black = truth, blue dashed = EKF estimate, green dashed = internal target or engagement band, red dashed = limit or tolerance, magenta dotted = supervisor stage switch. The numbers quoted are in \texttt{report\_summary\_<variant>.txt} and the output of \texttt{run\_all\_tests}. \texttt{README.md} in this folder lists which figure goes where in the report.

\section{Estimator figures}
Common setup: analytic truth trajectory (quintic staged approach from $(-2~\text{m}, 0.3~\text{m}, 15^\circ)$ to the standoff and on to the dock), simulated 200~Hz IMU with white noise and residual bias, 15~Hz fiducial fixes with range-dependent noise and a 50~ms delay, a 2~s marker dropout at $t = 18$ to $20$~s, handover initialisation (truth plus a draw from $\mathbf P_0$).

\entry{ekf_demo_plane.png}{Plane view in the target frame: truth path, EKF estimate, the fiducial fixes as points, the docked pose.}{\texttt{make\_ekf\_figures}, demo parameters.}{Context figure: the estimate rides on the truth, the fixes scatter with the range-dependent noise, the estimate coasts through the dropout.}

\entry{ekf_fullscale_plane.png}{As above, full-scale parameters (GQ7-class IMU, 5~mm fiducial noise at 1~m).}{\texttt{make\_ekf\_figures}, full scale.}{Same role for the full-scale instance.}

\entry{ekf_demo_errors.png}{Estimation error of $x$, $y$, $\theta$, $v_x$, $v_y$ against time with the filter's own $\pm 3\sigma$ bounds; grey band = dropout; green = internal target.}{Same run as the plane view.}{Consistency by inspection (errors inside the bounds about 99\,\% of the time), error growth during the dropout and re-convergence after it, and the 10 to 15~mm/s velocity-error level the MPC inherits through its initial increment.}

\entry{ekf_fullscale_errors.png}{As above, full scale.}{\texttt{make\_ekf\_figures}, full scale.}{Same role; note the smaller velocity error from the better IMU.}

\entry{ekf_demo_nees.png}{Top: single-run NEES against time with the 95\,\% chi-square bounds (5 dof). Bottom: mean NEES of 20 Monte Carlo runs with the band.}{20 runs with different noise seeds, same trajectory.}{The course's consistency test, used to choose the velocity process noise $q_v$: demo mean 5.32, inside the band.}

\entry{ekf_fullscale_nees.png}{As above, full scale.}{20 runs.}{Full-scale mean 3.21, slightly below the band: the filter is conservative at the chosen $q_v$, kept for margin against the IMU-bias assumption.}

\entry{ekf_demo_delay.png}{Position-error magnitude during the constant-velocity cruise with delay compensation on and off, with the line speed$\times\tau_d$; true speed in the lower panel.}{Same sensor record processed twice, with and without the buffered innovation.}{Justifies the ring-buffer delayed innovation: 10~mm lag bias at 0.2~m/s without it, comparable to the demo tolerance; 1.6~mm with it.}

\entry{ekf_fullscale_delay.png}{As above, full scale.}{Same.}{Same conclusion at full scale.}

\entry{ekf_demo_dropout.png}{Zoom on the 2~s dropout: position-error magnitude against the $3\sigma$ bound; $v_x$ error below.}{Same run as the errors figure.}{The demo leaves its 12~mm tolerance after about 0.2~s of coasting, so the demonstrator needs continuous marker visibility and the supervisor holds early.}

\entry{ekf_fullscale_dropout.png}{As above, full scale.}{Same.}{Full scale stays inside the 80~mm tolerance for the whole 2~s, which sets the supervisor's abort time.}

\entry{ekf_demo_target_motion.png}{The target starts moving at 0.2~m/s at $t_{\mathrm{step}}$, which the robot's IMU cannot see: $x$ error and $v_x$ error over the next 6~s.}{\texttt{apply\_target\_velocity\_step} on the truth trajectory, same sensors.}{The stationary-target assumption degrades gracefully: the filter re-converges from fixes alone within 1.2~s, no fix rejected. Also the evidence for the gate value: a tight gate rejected every fix here and the filter diverged.}

\entry{ekf_fullscale_target_motion.png}{As above, full scale.}{Same.}{Same conclusion.}

\section{Controller figures}
Common setup: closed loop with the EKF in the loop from $(-2~\text{m}, 0.3~\text{m}, 15^\circ)$, seed 1, truth plant with first-order wheel-speed loops (0.3~s full scale, 0.1~s demo) and a slip offset, three-stage supervisor.

\entry{mpc_demo_plane.png}{Path in the target frame: true path, EKF estimate, and the MPC's predicted trajectories (red) at ticks 1, 40, 120, 200 and 300; standoff and docked markers.}{\texttt{make\_mpc\_figures}, demo.}{Receding horizon with successive linearisation: each red curve is the linear model's 2~s forecast from that tick; the forecasts curve as the linearisation point moves and converge onto the path.}

\entry{mpc_fullscale_plane.png}{As above, full scale.}{Same.}{Same role.}

\entry{mpc_demo_timeseries.png}{$x$, $y$, $\theta$ against time, true and estimated, with the stage switches, the standoff and contact lines and the internal targets.}{Same run.}{The two-stage approach: alignment at the standoff, final approach at 50~mm/s, docking at 12.4~s. The heading settles within $\pm 1^\circ$ after the initial turn.}

\entry{mpc_fullscale_timeseries.png}{As above, full scale; docking at 14.5~s.}{Same.}{The $\pm 1^\circ$ heading ripple is the yaw loop against the 0.3~s wheel lag the kinematic model omits; the lag model of the design document removes it.}

\entry{mpc_demo_commands.png}{Body-velocity commands with the approach and final caps; increments with the translation and yaw limits; the four wheel speeds with the wheel limit (the course's output / input / increment layout).}{Same run, sampled at the controller ticks.}{Which constraints are active when: the increment limit during the initial turn and the deceleration at the standoff; the L1 speed cap during the approach. The wheel limit is never exceeded.}

\entry{mpc_fullscale_commands.png}{As above, full scale.}{Same.}{At full scale the wheel limit (4~rad/s) is the active constraint during the approach, which is where the inverse kinematics inside the controller matters.}

\entry{mpc_demo_rw.png}{$x(t)$ and the forward command for $r_w = 0.1$, 1 and 10, with the docking times.}{Three closed-loop runs, same seed.}{The course's single tuning knob: smaller $r_w$ gives more aggressive increments; the design uses $r_w = 1$.}

\entry{mpc_fullscale_rw.png}{As above, full scale.}{Same.}{Same role.}

\entry{mpc_demo_tick.png}{One tick's predicted $x$ over the 2~s horizon (red) against what actually happened (black), with the tick and the end of the control horizon marked.}{Prediction stored at tick 200 of the nominal run.}{Prediction quality of the per-tick linear model, the Example~5.3 flavour of the notes.}

\entry{mpc_fullscale_tick.png}{As above, full scale.}{Same.}{Same role.}

\section{Report figures}
All closed loop with the EKF in the loop, produced by \texttt{make\_report\_figures}.

\entry{report_demo_montecarlo.png}{30 runs from start poses drawn uniformly in $x\in[-2.5,-0.8]$~m, $y\in[-0.4,0.4]$~m, $\theta\in[-20^\circ,20^\circ]$ with independent noise seeds. Top left: final lateral and yaw error against the coupler tolerance (red box) and the internal target (green box); then docking time, penetration past first contact (red line: end of the 10~mm engagement band) and closing speed at contact (red line: final cap).}{\texttt{make\_report\_figures}, demo.}{The main repeatability result for the demo: 30/30 inside the tolerance, 26/30 inside the 4~mm target; lateral 2.4~mm rms, 5.6~mm max; yaw $\le 0.26^\circ$; docking 12.0~s mean, 15.2~s max; penetration $\le 10.1$~mm; contact $\le 43$~mm/s; no solver fallback. The four runs outside the target are estimator-limited: the assumed 4~mm fiducial noise per fix equals the target.}

\entry{report_fullscale_montecarlo.png}{As above, full scale.}{Same.}{30/30 inside tolerance and target; lateral 3.5~mm rms, 9.2~mm max against the 27~mm target; yaw $\le 0.61^\circ$; docking 13.3~s mean, 17.9~s max; penetration $\le 8.2$~mm; contact $\le 12$~mm/s.}

\entry{report_demo_sensitivity.png}{$3\times 4$ grid: lateral error at docking, penetration and docking time against the wheel-loop time constant $\tau$, the fiducial delay, the sensor-noise scale and the IMU-bias scale, one at a time, with the filter and controller unaware of the change. Blue dotted = nominal, green = target or band, red $\times$ = did not dock within 40~s.}{15 closed-loop runs from the nominal start pose, seed 21.}{The demo's margins: docks with $\tau$ up to 0.3~s, delay up to 0.1~s, noise up to $4\times$, bias up to $2\times$; does not dock at $\tau = 0.5$~s, delay 0.2~s or bias $4\times$. Hence the bench requirements: $\tau < 0.2$~s, latency $< 0.1$~s, residual bias within $2\times$ of 0.03~m/s$^2$.}

\entry{report_fullscale_sensitivity.png}{As above, full scale.}{Same.}{Full scale is insensitive over the ranges tested: $\tau$ up to 0.5~s, delay up to 0.2~s and bias up to $4\times$ all dock inside the targets; only noise at $4\times$ prevents the alignment condition being met within 40~s.}

\entry{report_demo_constraints.png}{The nominal run with the constraints in the optimiser (black) against the same controller with every limit removed from the QP and the command clipped to the limits afterwards (red): $x(t)$, $x$ near the contact plane, largest wheel speed.}{Two closed-loop runs, same seed; the clipped version uses the \texttt{unconstrained} option of \texttt{run\_closed\_loop}.}{An honest comparison: in the nominal run the two behave alike (penetration 3.9 against 4.2~mm) because the active limits are enforced equally by clipping and the output constraints are inactive. The constraints matter for the output rows, which clipping cannot enforce, and for the feasibility guarantee; the nominal performance does not rest on them.}

\entry{report_fullscale_constraints.png}{As above, full scale.}{Same.}{No penetration in either case; the clipped run is in fact 1~s faster. Same conclusion.}

\entry{report_demo_estimator.png}{Perfect state feedback (black) against the EKF in the loop, true pose (blue) and estimate (dotted): $x$, $y$, $\theta$.}{Two closed-loop runs, same seed, \texttt{use\_ekf} on and off.}{The estimator's cost: 15.3 against 13.5~s docking time and 1.3 against 0.2~mm final lateral error; the estimate stays within its covariance throughout. Certainty equivalence holds for this loop.}

\entry{report_fullscale_estimator.png}{As above, full scale.}{Same.}{15.8 against 14.6~s and 2.6 against 0.6~mm.}

\end{document}
